\documentclass[10pt,aspectratio=169]{beamer}
\usepackage[portuguese]{babel}
\usepackage{fontspec,unicode-math}
\setmainfont{Noto Sans}
\setsansfont{Noto Sans}
\setmathfont{STIX Math}
\usepackage{amsmath}
\usepackage{xcolor}
\definecolor{InsperRed}{HTML}{a01924}
\definecolor{Ink}{HTML}{232529}
\setbeamercolor{frametitle}{fg=Ink}
\setbeamercolor{structure}{fg=InsperRed}
\setbeamercolor{normal text}{fg=Ink,bg=white}
\setbeamerfont{frametitle}{size=\large,series=\bfseries}
\setbeamerfont{title}{size=\LARGE,series=\bfseries}
\setbeamertemplate{navigation symbols}{}
\setbeamertemplate{itemize item}{\color{InsperRed}$\bullet$}
\setbeamertemplate{footline}{%
  \hbox{\begin{beamercolorbox}[wd=.84\paperwidth,ht=2.7ex,dp=1.2ex,leftskip=.7cm]{normal text}\scriptsize Inteligência Artificial para Decisão · Insper · 2027-1\end{beamercolorbox}%
  \begin{beamercolorbox}[wd=.16\paperwidth,ht=2.7ex,dp=1.2ex,rightskip=.7cm plus1fil]{normal text}\hfill\scriptsize\insertframenumber\, /\, \inserttotalframenumber\end{beamercolorbox}}}
\setbeamersize{text margin left=12mm,text margin right=12mm}
\begin{document}
\begin{frame}[t]{}
\vspace{14mm}\textcolor{InsperRed}{\bfseries AULA 10}\par\vspace{4mm}{\LARGE\bfseries Avaliação experimental de aplicações de IA}\par\vspace{6mm}{\large Operações empresariais}\par\vspace{12mm}{\color{gray}Paulo C. Marques F. · 2027-1}
\end{frame}

\begin{frame}[t]{Objetivos e verificação}
\small
\begin{itemize}
\item Planejar comparações controladas, registrar ganhos e regressões, e avaliar respostas fundamentadas, custo e latência.
\item Resolver manualmente os exemplos do capítulo, identificando hipóteses e consequências.
\item Especificar a aplicação no Codex gráfico e confrontar ao menos um resultado com evidências.
\end{itemize}
\end{frame}

\begin{frame}[t]{Organização do encontro}
\small
\begin{itemize}
\item Apresentação dos fundamentos e resolução de problemas quantitativos no papel.
\item Laboratório supervisionado e testes produzidos por prompts na interface gráfica do Codex, sem terminal.
\item Atividade individual com dados diferentes e depósito no GitHub antes do quiz seguinte.
\end{itemize}
\end{frame}

\begin{frame}[t]{Comparações controladas}
\small
\begin{itemize}
\item Melhorar uma aplicação exige definir o que será mantido constante, o que será modificado e qual métrica representará o objetivo.
\item Uma avaliação antes/depois usando conjuntos de teste diferentes não identifica o efeito da mudança.
\item O desenho mínimo usa uma configuração de referência (baseline) e uma alternativa, ambas aplicadas às mesmas entradas.
\item Para uma métrica calculada em exemplos, escrevemos a diferença emparelhada .
\end{itemize}
\end{frame}

\begin{frame}[t]{Comparações controladas — continuação}
\small
\begin{itemize}
\item Ao comparar resultados individuais, o vetor ajuda a examinar casos em que a nova configuração melhorou ou piorou.
\item Uma média positiva pode esconder perdas graves em perguntas de risco elevado.
\end{itemize}
\end{frame}

\begin{frame}[t]{Separação de desenvolvimento e teste}
\small
\begin{itemize}
\item É inadequado testar dezenas de prompts no mesmo conjunto e apresentar o maior escore como resultado independente: a escolha foi feita usando o teste.
\item O desenho deve ter casos para especificação, casos para desenvolvimento e um conjunto final reservado.
\item O mesmo cuidado se aplica à seleção de tamanho de chunk, top-k e limiares.
\item Um avaliador por LLM (LLM-as-a-Judge) pode atribuir notas ou comparar respostas, mas sua saída é uma medição potencialmente enviesada e variável.
\end{itemize}
\end{frame}

\begin{frame}[t]{Separação de desenvolvimento e teste — continuação}
\small
\begin{itemize}
\item Devemos registrar o critério, fornecer evidência aos avaliadores, analisar discordâncias e manter casos de revisão humana.
\item O juiz não deve ser a única fonte do gabarito.
\end{itemize}
\end{frame}

\begin{frame}[t]{Resultados e trade-offs}
\small
\begin{itemize}
\item Uma versão que passa de 80\% para 84\% de acerto melhora quatro pontos percentuais, não quatro por cento relativos: a melhoria relativa é .
\item A clareza sobre denominadores é essencial.
\item Em sistemas reais, desempenho também envolve tempo de resposta, custo de execução, taxa de abstenção e consumo de contexto.
\item Não é correto comparar apenas qualidade sem considerar viabilidade operacional.
\end{itemize}
\end{frame}

\begin{frame}[t]{Diagnóstico por etapas}
\small
\begin{itemize}
\item Para um RAG de atendimento, se a política correta não foi recuperada, testar o índice e a consulta; se foi recuperada mas a resposta contradiz a política, investigar geração e instruções; se a resposta está correta mas viola…
\item O experimento deve registrar qual etapa mudou e qual etapa falhou.
\item Assim, a aula transforma a ideia vaga de ``melhorar o chatbot'' em um procedimento de avaliação reprodutível.
\end{itemize}
\end{frame}

\begin{frame}[t]{Exercício para resolver no papel 1}
\small
\small\textbf{10.1 Melhoria líquida e regressões}\par\vspace{4mm}Em 40 atendimentos, a configuração A acerta 26. B corrige oito falhas de A, mas erra três casos antes corretos. Calcule acertos de B, taxas de A e B e ganho em pontos percentuais. Por que olhar somente a taxa de B esconde um risco operacional?
\end{frame}

\begin{frame}[t]{Exercício para resolver no papel 2}
\small
\small\textbf{10.2 Ganho absoluto e relativo}\par\vspace{4mm}A taxa de respostas fundamentadas cresce de 60\% para 72\%. Calcule (a) ganho em pontos percentuais; (b) crescimento relativo; (c) a quantidade de casos adicionais em 50 solicitações, admitindo as proporções como frequências observadas.
\end{frame}

\begin{frame}[t]{Laboratório e caso aplicado}
\small
\begin{itemize}
\item Área de aplicação: Operações empresariais.
\item Definir previamente desfecho primário, rubrica de citação e limite operacional de latência.
\item Solicitar ao Codex comparação pareada, matriz de transições acerto–erro e relatórios de melhorias/regressões por ID.
\item Inspecionar casos de documento desatualizado, regra de exceção e afirmação elegante sem evidência.
\end{itemize}
\end{frame}

\begin{frame}[t]{Prompt de desenvolvimento — exemplo}
\small
\small\textbf{Instrução fornecida ao Codex:}\par\vspace{5mm}Compare duas configurações de triagem sobre os mesmos 40 casos. Calcule métricas pareadas, melhorias, regressões, fundamentação e tempos registrados, com relatórios por ID. Use previsões congeladas para que o experimento funcione sem chamada de modelo; não altere simultaneamente os dados.
\end{frame}

\begin{frame}[t]{Critérios de aceitação e validação}
\small
\begin{itemize}
\item Mesmos IDs e rótulos nas duas condições; regressões identificadas individualmente; latência descrita por mediana e quantis, se fornecidos; nenhuma inferência causal atribuída a mudanças múltiplas.
\item Inspecionar argumentos, saídas e rastros das ferramentas quando houver chamadas externas.
\item Comparar uma saída com a resolução manual do exemplo conceitual.
\end{itemize}
\end{frame}

\begin{frame}[t]{Atividade individual em casa}
\small
\begin{itemize}
\item Novo conjunto com 30 solicitações e regras fictícias de assistência técnica; saídas pareadas A/B e tempos medidos fornecidos pelo professor.
\item Resolver as questões de interpretação, com cálculos e evidência documental.
\item Depositar no GitHub os arquivos solicitados; a entrega habilita o quiz, sem valer nota.
\end{itemize}
\end{frame}

\begin{frame}[t]{Preparação para a avaliação escrita}
\small
\begin{itemize}
\item Exercícios do caderno: 10.1 Melhoria líquida e regressões; 10.2 Ganho absoluto e relativo; 10.3 Duas mudanças simultâneas.
\item Formalizar o raciocínio e interpretar resultados; não será exigida escrita manual de código.
\item No quiz seguinte serão avaliados conceitos desta aula e da atividade feita em casa.
\end{itemize}
\end{frame}

\end{document}
